\section*{[2024 EXAM] Problem 2: Cubic Lagrange finite element in 2D}
In this question we focus on the definition and construction of finite elements.
\subsection*{(a) Ciarlet's definition of a finite element}
\textbf{Ciarlet finite element}
Let
\begin{itemize}
    \item $K \subset \mathbb{R}^n$ be a closed, bounded set with nonempty interior and
          piecewise smooth boundary;
    \item $P$ be a finite-dimensional space of real-valued functions on $K$;
    \item $N = \{N_1,\dots,N_m\}$ be a set of linear functionals $N_i : P \to \mathbb{R}$ such
          that $\{N_i\}$ form a basis of the dual space $P^\ast$.
\end{itemize}
Then the triple $(K,P,N)$ is called a \emph{finite element} (a Ciarlet triple).
Equivalently, $N$ \emph{determines} $P$ in the sense that
\[
    p\in P,\quad N_i(p) = 0 \ \forall i \quad \Rightarrow\quad p \equiv 0.
\]
\subsection*{(b) Cubic Lagrange finite element in 2D}
We describe the cubic Lagrange finite element over a triangle $K \subset \mathbb{R}^2$.
\textbf{Geometry.}
Let $K$ be a nondegenerate triangle with vertices $a_1,a_2,a_3$.
Let $e_i$ be the edge opposite $a_i$.
\textbf{Polynomial space.}
We take
\[
    P := \mathbb{P}_3(K)
    = \{ \text{polynomials on }K \text{ of total degree }\le 3 \}.
\]
In 2D,
\[
    \dim \mathbb{P}_3 = \frac{(3+1)(3+2)}{2} = 10.
\]
\textbf{Degrees of freedom (nodal functionals).}
We use point evaluations at $10$ nodal points:
\begin{itemize}
    \item \emph{Vertices} (3 DoFs):
          \[
              N_i(v) := v(a_i), \quad i=1,2,3.
          \]
    \item \emph{Edge nodes} (6 DoFs):
          On each edge $e_i$ we place two distinct interior points, for instance at
          barycentric coordinates $(2/3,1/3,0)$ and $(1/3,2/3,0)$ relative to
          the pair of vertices defining the edge. Denote these by $m_{i1}$ and $m_{i2}$.
          Then
          \[
              N_{i,j}(v) := v(m_{ij}), \quad i=1,2,3,\ j=1,2.
          \]
    \item \emph{Interior node} (1 DoF):
          Let $b$ be the barycenter
          \[
              b := \frac{1}{3}(a_1 + a_2 + a_3).
          \]
          Define
          \[
              N_b(v) := v(b).
          \]
\end{itemize}
Altogether we have $3 + 6 + 1 = 10$ nodal functionals, matching $\dim P$.
\emph{Pictographically:} one draws a triangle with
\begin{itemize}
    \item a node at each vertex,
    \item two nodes equally spaced along each edge (at $1/3$ and $2/3$ of the edge length),
    \item one node in the interior (the barycenter).
\end{itemize}
Every degree of freedom is the value of $v$ at one of these nodes.
Thus the cubic Lagrange finite element in 2D is the Ciarlet triple
\[
    (K, \mathbb{P}_3(K), N),
\]
where $N$ consists of the 10 point-evaluation functionals described above.
\subsection*{(c) Cubic Lagrange elements satisfy Ciarlet's definition}
We must show that the triple $(K,P,N)$ from (b) is a finite element in the sense of Ciarlet, i.e.
\begin{itemize}
    \item $K$ is an admissible element domain;
    \item $P=\mathbb{P}_3(K)$ is finite-dimensional;
    \item $N$ is a basis of $P^\ast$ (equivalently: $N$ determines $P$).
\end{itemize}
The first two items are immediate, so the key point is:
\[
    \boxed{
        p \in \mathbb{P}_3(K),\ N(p) = 0\ \forall N\in N
        \quad \Rightarrow\quad
        p \equiv 0.
    }
\]
\textbf{Setup.}
Let $p \in \mathbb{P}_3(K)$ satisfy
\[
    p(a_i) = 0,\quad
    p(m_{ij}) = 0,\quad
    p(b) = 0,
\]
for all vertices $a_i$, edge nodes $m_{ij}$, and the barycenter $b$.
Let $e_1,e_2,e_3$ be the three edges of $K$, and let $L_i(x)=0$ be an affine equation for
the line supporting the edge $e_i$ (a hyperplane in $\mathbb{R}^2$).
We will use the given:
\begin{itemize}
    \item \emph{Definition:} A hyperplane is a set of the form $\{x : L(x)=0\}$ for non-degenerate
          linear $L$.
    \item \emph{Deconstruction lemma:} If $P$ is a polynomial of degree $d\ge 1$ vanishing on the
          hyperplane $L(x)=0$, then $P = L Q$ with $\deg Q = d-1$.
\end{itemize}
\textbf{Step 1: $p$ vanishes on edge $e_1$ and factors as $p = L_1 q_1$.}
Restrict $p$ to the edge $e_1$ (between vertices $a_2$ and $a_3$).
On this edge we have $4$ distinct nodal points:
\[
    a_2,\ a_3,\ m_{1,1},\ m_{1,2}.
\]
The restriction $p|_{e_1}$ is a univariate polynomial of degree at most $3$.
Since $p$ vanishes at four distinct points on $e_1$, it must vanish identically on $e_1$:
\[
    p(x) = 0 \quad \forall x \in e_1.
\]
Hence $p$ vanishes on the hyperplane $L_1(x)=0$.
By the deconstruction lemma,
\[
    p = L_1 q_1,
\]
for some polynomial $q_1$ of degree at most $2$.

\textbf{Step 2: $q_1$ vanishes on edge $e_2$ and factors as $q_1 = L_2 q_2$.}
Now restrict $p$ to edge $e_2$. The nodal points on $e_2$ are
\[
    a_1,\ a_3,\ m_{2,1},\ m_{2,2}.
\]
We know $p$ vanishes at all of these.
Along $e_2$ we have $L_2(x)=0$, while $L_1(x)=0$ only at the endpoint where $e_1$ and $e_2$
meet (one vertex). At the other three nodal points on $e_2$ (one vertex and two interior
edge points), $L_1\neq 0$.
For those three points $x$ we have
\[
    0 = p(x) = L_1(x)\,q_1(x), \quad L_1(x) \neq 0 \quad\Rightarrow\quad q_1(x)=0.
\]
Thus $q_1$ (a polynomial of degree $\le 2$) has three distinct zeros on the line $e_2$, so
$q_1$ vanishes identically on $e_2$.
Hence $q_1$ vanishes on hyperplane $L_2(x)=0$. By the deconstruction lemma again,
\[
    q_1 = L_2 q_2,
\]
for some polynomial $q_2$ of degree at most $1$.
Combining Steps 1 and 2:
\[
    p = L_1 L_2 q_2.
\]
\textbf{Step 3: $q_2$ vanishes on edge $e_3$ and factors as $q_2 = L_3 c$.}
Consider edge $e_3$. Its four nodal points are
\[
    a_1,\ a_2,\ m_{3,1},\ m_{3,2}.
\]
We know $p$ vanishes at all of these:
\[
    p(x) = L_1(x)L_2(x)q_2(x) = 0.
\]
At the two interior edge nodes $m_{3,1},m_{3,2}$, neither $L_1$ nor $L_2$ vanishes
(these points are not on the other two edges). Hence
\[
    L_1(m_{3,j})\neq 0,\quad L_2(m_{3,j})\neq 0,\quad p(m_{3,j})=0
    \quad \Rightarrow\quad q_2(m_{3,j}) = 0 \quad (j=1,2).
\]
Thus $q_2$ (degree $\le 1$) has two distinct zeros on the line $e_3$, so its restriction
to $e_3$ vanishes identically. Hence $q_2$ vanishes on the hyperplane $L_3(x)=0$, and the
deconstruction lemma gives
\[
    q_2 = L_3 c,
\]
for some constant $c\in\mathbb{R}$.
Therefore
\[
    p = L_1 L_2 L_3 c.
\]
\textbf{Step 4: Use the interior node to conclude $c=0$ and hence $p\equiv 0$.}
Finally, evaluate $p$ at the barycenter $b$. This point lies strictly inside $K$, so it
is not on any of the edges; in particular,
\[
    L_1(b)\neq 0,\quad L_2(b)\neq 0,\quad L_3(b)\neq 0.
\]
We know $p(b)=0$, so
\[
    0 = p(b) = c\,L_1(b)L_2(b)L_3(b).
\]
Since $L_1(b)L_2(b)L_3(b)\neq 0$, it must be that $c=0$. Thus
\[
    p \equiv 0 \quad\text{on }K.
\]
\textbf{Conclusion.}
We have shown that
\[
    p\in\mathbb{P}_3(K),\quad N(p)=0\ \forall N\in N
    \quad \Rightarrow\quad p\equiv 0.
\]
Therefore $N$ determines $P$. Since $\#N = 10 = \dim P$, the functionals in $N$ are
linearly independent and form a basis of $P^\ast$.
Hence Ciarlet's definition is satisfied and the cubic Lagrange element
$(K,\mathbb{P}_3(K),N)$ is a valid finite element.
\subsection*{(d) Mesh-level degrees of freedom for a conforming $H^1_0$ space}
We now consider a small triangulation (mesh) of $\Omega$ with at least four triangular
elements and place a cubic Lagrange element on each triangle.
\textbf{Local degrees of freedom per element} (from (b)):
\begin{itemize}
    \item 3 vertex nodes,
    \item 6 edge nodes (2 on each edge),
    \item 1 interior node.
\end{itemize}
\textbf{Conforming subspace of $H^1_0(\Omega)$.}
To obtain a conforming finite element space $V_h \subset H^1_0(\Omega)$, the finite element functions must be \emph{continuous} across element interfaces and vanish on the boundary $\partial\Omega$.
This translates into how the degrees of freedom are treated on the mesh:
\begin{itemize}
    \item \emph{Interior element nodes (barycenter of each triangle):}
          These lie strictly inside each element, so they are \emph{not} shared with neighbouring elements. Each such interior node corresponds to its own global degree of freedom.
    \item \emph{Vertices in the interior of $\Omega$:}
          Mesh vertices that are shared by several triangles correspond to \emph{one} global node.
          The nodal value $u_h(a)$ at such a vertex $a$ is shared between all elements touching $a$. This enforces continuity at interior vertices.
    \item \emph{Interior edge nodes (two nodes on each interior edge):}
          Each interior edge $e$ is shared by two triangles. The two cubic Lagrange nodes on $e$ are common to both elements.
          Hence each of these nodes corresponds to a single global degree of freedom, shared by the adjacent elements.
          This ensures continuity of $u_h$ along interior edges.
    \item \emph{Boundary nodes (on $\partial\Omega$):}
          All vertices and edge nodes lying on the boundary $\partial\Omega$ must satisfy the homogeneous Dirichlet condition $u=0$.
          Therefore, these boundary nodal values are \emph{fixed to zero} (not free DoFs).
\end{itemize}
\textbf{Answer in words (for a sketch of 4 elements):}
\begin{itemize}
    \item Draw a small patch of at least four triangles.
    \item Mark all cubic Lagrange nodes (3 vertices, 6 edge points, 1 interior point per element).
    \item \emph{Shared DoFs (to enforce continuity):}
          \begin{itemize}
              \item all vertex nodes in the interior of the domain,
              \item all edge nodes on interior edges (each such node is shared by the two adjacent
                    elements).
          \end{itemize}
    \item \emph{Local-only DoFs:}
          \begin{itemize}
              \item the interior node (barycenter) of each triangle (belongs to that element only).
          \end{itemize}
    \item \emph{Fixed DoFs (for $H^1_0$):}
          \begin{itemize}
              \item all nodes lying on $\partial\Omega$ (vertices and boundary edge nodes) are fixed to
                    value $0$.
          \end{itemize}
\end{itemize}
With these identifications, the global finite element space $V_h$ consists of functions
that are continuous across element boundaries and vanish on $\partial\Omega$, hence
\[
    V_h \subset H^1_0(\Omega),
\]
and the cubic Lagrange finite element space is conforming.
